% Generated draft: requires mathematical and source review.
\documentclass[11pt]{article}
\usepackage[margin=25mm]{geometry}
\usepackage{fontspec}
\setmainfont{Latin Modern Roman}
\usepackage{amsmath,amssymb,mathtools,unicode-math}
\setmathfont{Latin Modern Math}
\usepackage{enumitem,needspace,xurl,hyperref}
\setlength{\parindent}{0pt}
\setlength{\parskip}{5pt}
\setlength{\emergencystretch}{3em}
\newcommand{\N}{\mathbb{N}}
\newcommand{\Z}{\mathbb{Z}}
\newcommand{\Q}{\mathbb{Q}}
\newcommand{\R}{\mathbb{R}}
\newcommand{\C}{\mathbb{C}}
\newcommand{\F}{\mathbb{F}}
\newcommand{\A}{\mathbb{A}}
\newcommand{\PP}{\mathbb{P}}
\newcommand{\E}{\mathbb{E}}
\newcommand{\eps}{\varepsilon}
\newcommand{\dd}{\,\mathrm{d}}
\begin{document}
\section*{UnsolvedMath ID: 30003851}
\textbf{Source problem code:} catalogue entry
\subsection{Definitions and mathematical statement}
Let \(\operatorname{Homeo}^{+}(\R)\) be the group of increasing homeomorphisms of \(\R\), with composition as its operation. Define \(\alpha,\beta\in\operatorname{Homeo}^{+}(\R)\) by
\[
\alpha(t)=t+\frac12,\qquad
\beta(t)=
\begin{cases}
t,&t\leq0,\\
\dfrac{t}{1-t},&0\leq t\leq\frac12,\\
3-\dfrac1t,&\frac12\leq t\leq1,\\
t+1,&t\geq1.
\end{cases}
\]
The formulas for \(\beta\) agree at the shared endpoints \(0,\frac12,1\). Let
\[
G=\langle\alpha,\beta\rangle\leq\operatorname{Homeo}^{+}(\R).
\]
Regard \(G\) as a discrete group. It is amenable if there exists a finitely additive function \(\mu:\mathcal P(G)\to[0,1]\) such that \(\mu(G)=1\) and \(\mu(gA)=\mu(A)\) for every \(g\in G\) and every subset \(A\subseteq G\), where \(gA=\{ga:a\in A\}\).

\textbf{Precise yes-or-no question.} Is \(G\) nonamenable; that is, does no such \(\mu\) exist?

For the accompanying comparison, Thompson's group \(F\) denotes the group of increasing piecewise-linear homeomorphisms of \([0,1]\) with finitely many breakpoints, all dyadic rational, and slopes in \(\{2^k:k\in\Z\}\). A dyadic rational is a number \(m/2^n\), where \(m\in\Z\) and \(n\) is a nonnegative integer.

\textbf{Unresolved qualifications of the source question.} The record describes \(G\) as a subgroup of Monod's group, but supplies no definition of that ambient group. The explicit generators nevertheless determine \(G\) in \(\operatorname{Homeo}^{+}(\R)\). The original question also asks whether this is the simplest nonamenable example going slightly beyond \(F\), and why the coefficient \(1/2\) is important. No criterion of simplicity, precise relationship to \(F\), or comparison obtained by changing \(1/2\) is specified. Thus these accompanying requests remain qualitative questions requiring clarification; they do not assert a specified embedding, extension, or other relationship between \(F\) and \(G\). The number \(1/2\) occurs both as the translation amount in \(\alpha\) and as a breakpoint in \(\beta\); the record does not specify a parameterized replacement of these occurrences.
\subsection{Short English statement}
Determine whether the explicitly generated group \(G=\langle\alpha,\beta\rangle\), considered as a discrete group, is nonamenable. The source also asks why \(1/2\) matters and whether this is the simplest example slightly beyond Thompson's group \(F\), but supplies no precise comparison or simplicity criterion.
\subsection{Sources}
\begin{itemize}
\item Local UnsolvedMath registry, permanent ID \(30003851\), source problem code \texttt{catalogue entry}, title: ``A Nonamenable Extension Beyond Thompson's Group.'' The supplied record belongs to the Oberwolfach Reports Open Problems collection and gives proposed year 2018; no individual proposer is identified.
\item Supplied source citation: \emph{Cohomological and Metric Properties of Groups of Homeomorphisms of R} (2019), \emph{Oberwolfach Reports}, pp.\ 1579--1633, DOI 10.4171/owr/2018/26. \url{https://doi.org/10.4171/owr/2018/26}. The supplied background describes these reports as records of mathematical workshops and their open-problem discussions.
\item Supplied literature source: ``Oberwolfach Report 26/2018.'' \url{https://ems.press/content/serial-article-files/46748?nt=1}.
\item Literature assessment dated \texttt{2026-08-21}: the supplied curation labels the problem open and describes the questions about the simplest nonamenable example beyond \(F\) and the importance of \(1/2\) as unresolved. This is supplied dated curation, not an independently verified literature assessment. The record itself requests independent verification of time-sensitive status and cited claims.
\item The supplied statement carries the registry label \texttt{corrected\_verified}, with a claim of manual semantic, typographic, and LaTeX review. This records the supplied review claim; it does not establish mathematical correctness or independent source verification for this draft.
\end{itemize}
\end{document}
